\begin{thebibliography}{99}
\addcontentsline{toc}{section}{参考资料与实现依据}
\bibitem{course} 课程实验手册．第五--六章：基于视觉与AI定位的多机协同控制实验．仓库根目录同名\code{.docx}文件。
\bibitem{repo} 项目协作仓库．\url{https://github.com/yky666/intelligent_robot_system}．本轮整合基线：\code{139c0b7}；原分析图与快照另注明历史提交。
\bibitem{maze} TonyPi标记航路实现及实验状态．\code{humanoid/maze_route/README.md}，\code{marker_route.py}，\code{markers.py}，\code{robust_camera.py}，\code{tests/test_maze_route.py}。
\bibitem{follow2} 第一组地图跟随实现与部署记录．\code{turtlebot/follow2/multi_robot_following/}，\code{docs/vm_follow2.md}，\code{docs/turtlebot_debug_commands.md}。
\bibitem{follow1} AR直接跟随实验总结、独立报告与功能包．\code{docs/ar_marker_follow_experiment_summary.md}，\code{latex/follow1/ar_marker_report.tex}，\code{turtlebot/follow1/ar_follower/}与\code{ar_marker_following/}。前两份文档同时包含不同版本内容，本报告按源码区分。
\bibitem{follow3} RGB-D实时跟随方案、状态与验收要求．\code{turtlebot/follow3/README.md}及对应功能包。
\bibitem{gmapping} Grisetti G, Stachniss C, Burgard W. Improved Techniques for Grid Mapping with Rao-Blackwellized Particle Filters. IEEE Transactions on Robotics, 2007, 23(1):34--46.
\bibitem{dwa} Fox D, Burgard W, Thrun S. The Dynamic Window Approach to Collision Avoidance. IEEE Robotics and Automation Magazine, 1997, 4(1):23--33.
\bibitem{repo-current} 项目主分支复核．提交\code{43a7771baa2b33afc203b54543001995838879fd}，2026年10月17日读取．\url{https://github.com/yky666/intelligent_robot_system/tree/43a7771baa2b33afc203b54543001995838879fd}。
\bibitem{follow4} 方式四小组投稿：TurtleBot2视觉跟随与队列行驶．目录\code{turtlebot/follow4/team_turtlebot_follow_plus/}中的方式四README、\code{src/ar_laser_follower.cpp}与\code{launch/follow_ar.launch}；当前\code{sections/50_follow4.tex}含组员和实验记录。早期地图位姿适配方案仅作历史对照。
\bibitem{v1} 第一小组参考稿．\code{v1.pdf}，80页，封面日期2026年10月17日．本次由用户提供，用于第一小组杨开元、林楗棋、曹淇涵、刘小凡分工及技术内容参考，不代替其他小组报告。
\end{thebibliography}
